% Adapted to match the CIS 5511 PT_06 report format provided by the student.
% Compile with: pdflatex PT_07.tex (twice). Keep cisXXX.cls beside this file,
% if the course class is available.
\IfFileExists{cisXXX.cls}{
  \documentclass{cisXXX}
}{
  \documentclass[11pt]{article}
  \usepackage[margin=1in]{geometry}
  \newcommand{\HWauthor}[2]{\author{#1\\\texttt{#2}}}
  \newcommand{\HWno}[1]{\title{CIS 5511 --- Homework #1\\Order-$m$ B-Tree of Integers}}
  \newcommand{\HWcourse}[1]{}
}
\HWauthor{Md. Mizanur Rahman}{md.mizanur.rahman@temple.edu}
\HWno{7}
\HWcourse{CIS 5511}
\date{}
\usepackage{amsmath,amssymb,array,booktabs,listings,hyperref,tikz}
\hypersetup{colorlinks=true,urlcolor=blue,linkcolor=black}
\lstset{basicstyle=\ttfamily\footnotesize,breaklines=true,columns=fullflexible,
 keepspaces=true,showstringspaces=false,frame=single,tabsize=4}
\setlength{\emergencystretch}{3em}
\linespread{0.98}

\begin{document}
\maketitle
\section*{Order-$m$ B-Tree of Integers}

\subsection*{A. Design: representation and invariants}
An order-$m$ B-tree is implemented to store distinct integers, where the order
$m$ is specified when the tree is created. Following Lecture 09, $m$ denotes
the \emph{maximum number of children} of an internal node, not the minimum
degree. Each node therefore stores at most $m-1$ sorted keys. Except for the
root, every node stores at least $\lceil m/2\rceil-1$ keys. An internal node
containing $k$ keys has exactly $k+1$ ordered child references. All leaves
occur at the same depth, and each child's keys fall strictly between the
appropriate adjacent separator keys in its parent. The root is allowed to have
fewer keys than a non-root node. An empty B-tree is represented by an empty
root leaf.

\paragraph{Node representation.}
The class \texttt{BTreeNode} contains the following variables and method:
\begin{center}
\begin{tabular}{@{}p{0.27\linewidth}p{0.64\linewidth}@{}}
\toprule
Variable / method & Meaning \\
\midrule
\texttt{keys} & Sorted list of integer keys in the node.\\
\texttt{children} & Ordered list of child-node references; empty for a leaf.\\
\texttt{is\_leaf()} & Returns whether the node has no children.\\
\bottomrule
\end{tabular}
\end{center}

\paragraph{B-tree representation.}
The class \texttt{BTree} contains these variables:
\begin{center}
\begin{tabular}{@{}p{0.27\linewidth}p{0.64\linewidth}@{}}
\toprule
Variable & Meaning \\
\midrule
\texttt{m} & Order specified by \texttt{BTree(m)}; an integer at least 3.\\
\texttt{max\_keys} & Maximum number of keys per node: $m-1$.\\
\texttt{min\_keys} & Minimum number of keys in a non-root node: $\lceil m/2\rceil-1$.\\
\texttt{root} & Reference to the current root node.\\
\bottomrule
\end{tabular}
\end{center}

The required orders have the following occupancy limits:
\begin{center}
\begin{tabular}{@{}cccc@{}}
\toprule
Order & Maximum children & Maximum keys & Minimum non-root keys\\
\midrule
$m=4$ & 4 & 3 & 1\\
$m=5$ & 5 & 4 & 2\\
\bottomrule
\end{tabular}
\end{center}

\subsection*{Algorithms and their interfaces}
The following method names describe the implementation. The assignment does
not require detailed pseudocode or an algorithmic complexity analysis.
\begin{center}
\begin{tabular}{@{}p{0.33\linewidth}p{0.58\linewidth}@{}}
\toprule
Method & Main purpose\\
\midrule
\texttt{search(value)} & Find whether the integer occurs in the B-tree.\\
\texttt{insert(value)} & Insert a new integer; ignore a duplicate.\\
\texttt{\_insert(node, value)} & Insert recursively and propagate any split upward.\\
\texttt{delete(value)} & Delete an existing integer and update the root if needed.\\
\texttt{\_delete(node, value)} & Remove a leaf key or use the in-order successor for an internal key.\\
\texttt{\_repair\_underflow(parent, i)} & Borrow from a sibling or merge nodes after a deletion.\\
\texttt{display()} & Print the complete structure level by level.\\
\texttt{inorder()} & Return all stored keys in increasing order.\\
\texttt{validate()} & Check occupancy, ordering, child counts, and equal leaf depths.\\
\bottomrule
\end{tabular}
\end{center}

\paragraph{Insertion.}
Insertion searches for the appropriate leaf, places the new integer there in
sorted order, and splits a node if it temporarily has $m$ keys. The promoted
separator is the key at zero-based index $\lfloor m/2\rfloor$; this selects
the upper middle key when $m$ is even. A split may propagate upward. If the
root splits, a new root is created and the height increases by one.

\paragraph{Deletion.}
If the key is in an internal node, its in-order successor is copied into that
position and the successor is removed from the corresponding leaf. If a node
falls below the minimum number of keys, the implementation first attempts
to borrow from its right sibling, then its left sibling. If neither sibling
has an extra key, the node is merged with a sibling through a separating
parent key. Underflow repair may propagate upward, and an empty internal
root is replaced by its only child. These tie-breaking choices are
deterministic; another correct implementation may display a different valid
B-tree.

\subsection*{B. Programming and required tests}
The executable implementation is \texttt{btree.py}. Run it using:
\begin{lstlisting}
python btree.py
\end{lstlisting}
The program automatically performs two independent tests. Each test creates
an empty \texttt{BTree(m)}, calls \texttt{insert(i)} for $i=1,2,\ldots,20$
one by one, displays the structure, calls \texttt{delete(i)} for
$i=2,4,\ldots,20$ one by one, and displays the structure again. The
sample test procedure and corresponding console output are submitted separately
in \texttt{btree\_output.txt}; they are not duplicated in this report.

\subsection*{Summary of the two tests}
For $m=4$, each non-root node must contain 1--3 keys. For $m=5$, it must
contain 2--4 keys. Both tests completed every required operation successfully.
The following table summarizes the displayed trees (the root is shown in
square brackets, and the level count includes the root level):
\begin{center}
\begin{tabular}{@{}ccccc@{}}
\toprule
Order & Insertions & Deletions & Root after insertion & Root after deletion\\
\midrule
$m=4$ & $1$--$20$ & $2,4,\ldots,20$ & $[9]$ (3 levels) & $[9]$ (3 levels)\\
$m=5$ & $1$--$20$ & $2,4,\ldots,20$ & $[9]$ (3 levels) & $[5,11]$ (2 levels)\\
\bottomrule
\end{tabular}
\end{center}
After deleting the ten even integers, both trees contain exactly
$\{1,3,5,7,9,11,13,15,17,19\}$. The complete level-by-level structures
after insertion and deletion are provided in the separate text file.

\subsection*{Validation and reproducibility}
After each insertion and deletion, \texttt{validate()} verifies that every
node satisfies the occupancy bound, all keys are ordered and within their
subtree ranges, each internal node has one more child than key, and all leaves
have equal depth. The in-order traversal is also compared with the expected
integers. Both required tests passed. Additionally, 48,000 randomized insertions
and 48,000 randomized deletions across orders 3 through 10 passed the same
invariant and set-membership checks.

\subsection*{Files and execution}
The submission consists of three files: \texttt{PT\_07.pdf}, the compiled
report describing the B-tree design and test results; \texttt{btree.py}, the
Python implementation; and \texttt{btree\_output.txt}, containing the sample
test procedure and corresponding outputs. This \LaTeX{} source is retained
for editing. The course file \texttt{cisXXX.cls} may be used if available;
otherwise the report uses the standard article class.

\subsection*{Reference}
P. Wang, \emph{CIS 5511 Programming Techniques, Lecture 09: B-trees and
Amortized Analysis}.\\
\url{https://cis.temple.edu/~pwang/5511-PT/Lecture/5511-09.htm}

\subsection*{AI Use Acknowledgment}
Generative AI (ChatGPT) assisted with algorithm development, Python
implementation, testing, and LaTeX preparation.

\end{document}
